\documentclass{ximera}

\title{Methode der beginsnelheden}
\author{Mila Vervoort}
\license{CC: 0}         % replace with an appropriate license, or set it in xmPreamble

\begin{document}
\begin{abstract}
    Methode der beginsnelheden
\end{abstract}
\maketitle
\label{xim:Methode der beginsnelheden}

\section*{Methode der beginsnelheden}
\subsubsection*{Algemeen principe}

De methode der beginsnelheden is een techniek waarmee men de parameters van de snelheidsvergelijking kan bepalen zonder de volledige reactiecurve te analyseren.\\
De beginsnelheid $R_{A0}$ wordt gedefinieerd als de reactiesnelheid op het tijdstip $t=0$, wanneer de concentratie nog gelijk is aan de beginconcentratie $C_{A0}$.\\
Het centrale idee van deze methode is dat de beginsnelheid op een karakteristieke manier afhangt van de beginconcentratie. Door meerdere experimenten uit te voeren met verschillende waarden van $C_{A0}$ kan deze afhankelijkheid worden onderzocht en kunnen de kinetische parameters worden bepaald.

\begin{procedure}
\textbf{Methode der beginsnelheden}
De methode wordt toegepast volgens de volgende stappen:

\begin{enumerate}
    \item Voer meerdere experimenten uit met verschillende beginconcentraties $C_{A0}$.

    \item Bepaal voor elk experiment de beginsnelheid $R_{A,0}$ door differentiatie van de opgemeten data en extrapolatie naar $t \to 0$.

    \item Veronderstel een snelheidsvergelijking van de vorm
    \[
    -R_A = kC_A^n.
    \]

    \item Herschrijf deze vergelijking zodat een lineair verband ontstaat:
    \[
    \ln(-R_{A0})=\ln(k)+n\ln(C_{A0}).
    \]

    \item Maak een grafiek van $\ln(-R_{A0})$ als functie van $\ln(C_{A0})$. Indien de veronderstelde snelheidsvergelijking correct is, bekomt men een rechte waarvan:
    \begin{itemize}
        \item de helling gelijk is aan de reactieorde $n$;
        \item het snijpunt met de verticale as gelijk is aan $\ln(k)$.
    \end{itemize}
\end{enumerate}

\end{procedure}
\text{ } \\

Deze methode vereist slechts gegevens uit het begin van de reactie en is daarom bijzonder geschikt voor reacties waarbij de omgekeerde reactie of nevenreacties later een invloed hebben op de gemeten reactiesnelheid.

\end{document}